\section{行动与落地建议}

\subsection{短期行动清单}

\begin{table}[H]
\centering
\begin{tabular}{p{3cm}p{6cm}}
\toprule
时间窗 & 可交付物 \\
\midrule
Week 1 & 完成 data review + skill segmentation；设立 prompt log template \\
Week 2 & Launch evaluation dashboard，记录 multi-metric drift；设定 incident grade-book \\
Week 3 & 触发一次 emergent drill：LLM plan + skill exec 与 governance review \\
\bottomrule
\end{tabular}
\caption{短期行动建议。}
\end{table}

\begin{importantbox}{行动落地要点}
\begin{itemize}
    \item Prompt log + LLN plan 必须挂上 version/temperature/expected skill；
    \item 每次 emergent drill 都要有 review notes，形成 lessons learned；
    \item Incident grade-book 记录 trigger、redirect action、后续 guardrail change。
\end{itemize}
\end{importantbox}

\subsection{长期研究议程}

推荐以 quarterly thesis 方式梳理长期方向：1)scale-law + hierarchical optimizer；2)LLM-driven skill discovery；3)attention interpretability for subgoal.

\begin{knowledgebox}{长期研究示例}
\begin{itemize}
    \item Project Alpha：探索 hierarchical attention 模式；
    \item Project Beta：把 scale law 曲线映射到 optimizer schedule；
    \item Project Gamma：构建 emergent dashboard，帮助 governance panel 决策。
\end{itemize}
\end{knowledgebox}

\subsection{Checklist automation}

把短期 action checklist 自动化，例如用 Github Actions 每周检查 prompt log 是否同步、用 dashboards 展示 multi-metric drift，减少重复手动验证，工程师可以把精力集中在 incident response 上。

\begin{importantbox}{自动化 checklist 提醒}
\begin{itemize}
    \item prompt 更新触发 \texttt{check\_prompt\_versions} workflow；
    \item drift alert 达到 threshold 触发 \texttt{incident-drill}；
    \item \texttt{summary.md} 每周生成 lessons 概览，供 governance panel 审阅。
\end{itemize}
\end{importantbox}

\subsection{本章小结}

行动建议部分把理论转化为 timeline：短期 drill 快速闭环，长期项目则把 emergent insights 打造成可验证的 research output。

